\documentclass[12pt, a4paper]{article}
\usepackage[utf8]{inputenc}
\usepackage[margin=1in]{geometry}
\usepackage{amsmath}
\usepackage{amssymb}
\usepackage{graphicx}
\usepackage{array}
\usepackage{parskip}
\usepackage{hanging}
\usepackage{setspace}


\title{Research Project: Franciscan Eco-Theology and Environmental Attitudes}
\author{
A term paper submitted for the\\
RM in Management and Applied Microeconomics\\ [1.0cm]
Juan Taborda (), \\
Lidiia Levytska (), \\
Giovanna C. Pantalena (), \\
Kamilla Ospanova (50266102) \\[1.0cm]
Supervised by: Paul Behler \\[1.0cm]}

\begin{document}
\maketitle

\vspace{1.5cm}

\begin{abstract}
Our project focuses on the effect of exposure to the main
Mendicant orders of the 12th and 15th century - the Franciscan
and Dominican - on the contemporary environmental attitudes.
We find that that higher Franciscan is associated with less
pro-environmental attitudes, while higher Dominican exposure
is associated with more pro-environmental attitudes.
\end{abstract}

\newpage
\tableofcontents
Historical Background\\
Literature Review\\
Data and Methodology\\
Results\\
Limitations and Further Research\\
Conclusion\\


\newpage
\end{document}